\documentclass[12pt,a4paper]{article}

\usepackage[utf8]{inputenc}
\usepackage[T1]{fontenc}
\usepackage[french]{babel}
\usepackage[margin=2.5cm]{geometry}
\usepackage{listings}
\usepackage{xcolor}
\usepackage{hyperref}
\usepackage{parskip}
\usepackage{booktabs}
\usepackage{fancyhdr}
\usepackage{lmodern}

% ── style des blocs de code ───────────────────────────────────────────────────
\definecolor{codebg}{HTML}{F5F5F5}
\definecolor{keyword}{HTML}{0000CC}
\definecolor{comment}{HTML}{007700}
\definecolor{string}{HTML}{CC0000}
\definecolor{linenum}{HTML}{888888}

\lstdefinestyle{cstyle}{
    language=C,
    backgroundcolor=\color{codebg},
    basicstyle=\ttfamily\small,
    keywordstyle=\color{keyword}\bfseries,
    commentstyle=\color{comment}\itshape,
    stringstyle=\color{string},
    numberstyle=\tiny\color{linenum},
    numbers=left,
    numbersep=8pt,
    stepnumber=1,
    breaklines=true,
    breakatwhitespace=false,
    tabsize=4,
    showstringspaces=false,
    frame=single,
    framesep=4pt,
    xleftmargin=16pt,
    xrightmargin=4pt,
    captionpos=b,
}

\lstset{style=cstyle}

% ── en-tête et pied de page ───────────────────────────────────────────────────
\pagestyle{fancy}
\fancyhf{}
\rhead{TP3 --- Programmation Système}
\lhead{Gestion des processus en C}
\cfoot{\thepage}

% ── document ──────────────────────────────────────────────────────────────────
\begin{document}

\begin{titlepage}
    \centering
    \vspace*{2cm}
    {\Huge\bfseries TP3 : Gestion des processus\par}
    \vspace{0.5cm}
    {\Large Programmation Système en C\par}
    \vspace{2cm}
    \rule{\linewidth}{0.4pt}\\[0.4cm]
    {\large Trois programmes explorant les processus, les signaux et le système de fichiers\par}
    \rule{\linewidth}{0.4pt}
    \vfill
    {\large Ilyass Bougati\par}
    \vspace{0.3cm}
    {\large \today\par}
\end{titlepage}

\tableofcontents
\newpage

% ── Introduction ──────────────────────────────────────────────────────────────
\section{Introduction}

Ce rapport présente trois programmes en C développés dans le cadre du
troisième travail pratique de programmation système. Ces programmes couvrent
trois aspects complémentaires de la programmation UNIX :

\begin{itemize}
    \item \textbf{myls} --- une réimplémentation partielle de la commande
          \texttt{ls}, mettant en \oe{}uvre les API de répertoires et
          \texttt{stat}.
    \item \textbf{oroboros} --- un lanceur/tueur de processus fondé sur une
          liste chaînée, utilisant \texttt{fork}, \texttt{execvp} et
          l'envoi de signaux.
    \item \textbf{spam} --- un outil de stress qui forke un très grand nombre
          de processus fils afin de saturer la table des processus du noyau.
\end{itemize}

Les trois programmes sont compilés avec \texttt{gcc -Wall -Wextra} via le
\texttt{Makefile} commun fourni dans le projet.

% ── myls ──────────────────────────────────────────────────────────────────────
\section{\texttt{myls} --- Listage de répertoire}

\subsection{Objectif}

\texttt{myls} est un clone simplifié de l'utilitaire standard \texttt{ls}.
Il liste le contenu d'un répertoire en ignorant les entrées cachées (celles
dont le nom commence par un point) et supporte deux options :

\begin{center}
\begin{tabular}{@{}ll@{}}
\toprule
Option & Effet \\
\midrule
\texttt{-l} & Format long : permissions, liens, propriétaire, groupe, taille, date \\
\texttt{-R} & Récursif : descend dans les sous-répertoires \\
\bottomrule
\end{tabular}
\end{center}

\subsection{Appels système et API utilisés}

\begin{description}
    \item[\texttt{opendir} / \texttt{readdir} / \texttt{closedir}]
        Permettent d'itérer sur les entrées d'un répertoire via le flux
        \texttt{DIR} et la structure \texttt{struct dirent}.
    \item[\texttt{lstat}]
        Récupère les métadonnées d'un fichier (\texttt{struct stat}) sans
        suivre les liens symboliques, ce qui est le comportement correct
        pour \texttt{ls -l}.
    \item[\texttt{getpwuid} / \texttt{getgrgid}]
        Traduisent les valeurs numériques UID et GID en noms lisibles
        depuis \texttt{/etc/passwd} et \texttt{/etc/group}.
    \item[\texttt{localtime} / \texttt{strftime}]
        Formatent l'horodatage de dernière modification selon le style
        standard \texttt{MMM JJ HH:MM}.
\end{description}

\subsection{Remarques de conception}

Le mode récursif (\texttt{-R}) collecte les chemins des sous-répertoires dans
un tableau alloué dynamiquement (via \texttt{realloc}) lors d'un premier
parcours, puis les traite après l'appel à \texttt{closedir}. Cela évite de
maintenir plusieurs flux de répertoires ouverts simultanément.

Les permissions sont décodées bit à bit à l'aide des macros \texttt{S\_I*}
définies dans \texttt{<sys/stat.h>}, produisant la chaîne de dix caractères
habituelle (ex.\ \texttt{drwxr-xr-x}).

\subsection{Utilisation}

\begin{lstlisting}[language=bash, style=cstyle]
./myls                  # lister le repertoire courant
./myls /etc             # lister /etc
./myls -l /tmp          # format long
./myls -R .             # recursif depuis le repertoire courant
./myls -lR /usr/local   # long + recursif
\end{lstlisting}

\subsection{Code source}

\lstinputlisting[caption={\texttt{myls.c} --- utilitaire de listage de répertoire}]{myls.c}

% ── oroboros ──────────────────────────────────────────────────────────────────
\section{\texttt{oroboros} --- Le serpent des processus}

\subsection{Objectif}

\texttt{oroboros} prend une liste de noms de programmes en argument, construit
une liste chaînée à raison d'un nœud par programme, fork un processus fils
pour chaque nœud (en lançant le programme correspondant), affiche l'état de la
liste, puis --- sur appui d'une touche --- envoie \texttt{SIGTERM} à chaque
fils, attend une seconde, puis envoie \texttt{SIGKILL} aux survivants avant de
récolter tous les fils et de libérer la liste.

Le nom fait référence à l'ancien symbole du serpent qui se mord la queue : le
programme qui crée la vie la détruit aussi.

\subsection{Structure de données}

\begin{lstlisting}[caption={Nœud de la liste chaînée \texttt{AppNode}}]
typedef struct AppNode {
    char           *name;
    pid_t           pid;     /* renseigne apres fork */
    struct AppNode *next;
} AppNode;
\end{lstlisting}

Chaque nœud stocke le nom du programme (alloué sur le tas via
\texttt{strdup}), le PID assigné après \texttt{fork}, et un pointeur vers le
nœud suivant.

\subsection{Cycle de vie des processus}

\begin{enumerate}
    \item \textbf{Construction} --- \texttt{list\_append} parcourt la liste
          jusqu'à la queue et y ajoute un nœud pour chaque entrée de
          \texttt{argv}.
    \item \textbf{Lancement} --- \texttt{launch} appelle \texttt{fork} ; le
          fils redirige stdout et stderr vers \texttt{/dev/null} (pour ne pas
          polluer le terminal) puis appelle \texttt{execvp}, qui le remplace
          par le programme cible.
    \item \textbf{Arrêt} --- \texttt{kill\_all} effectue deux passes :
          d'abord \texttt{SIGTERM} (arrêt gracieux), puis après une seconde
          de délai \texttt{SIGKILL} (arrêt inconditionnel). Une boucle
          \texttt{waitpid(-1, NULL, WNOHANG)} récolte ensuite les zombies
          éventuels.
    \item \textbf{Libération} --- \texttt{list\_free} parcourt la liste et
          libère la chaîne de nom ainsi que le nœud lui-même.
\end{enumerate}

\subsection{Utilisation}

\begin{lstlisting}[language=bash, style=cstyle]
./oroboros top vim zsh nano
\end{lstlisting}

\subsection{Code source}

\lstinputlisting[caption={\texttt{oroboros.c} --- lanceur/tueur de processus par liste chaînée}]{oroboros.c}

% ── spam ──────────────────────────────────────────────────────────────────────
\section{\texttt{spam} --- Inondation de processus}

\subsection{Objectif}

\texttt{spam} est un outil de stress. Pour chaque nom de programme passé en
argument, il fork jusqu'à \texttt{COPIES} (1\,000\,000) processus fils,
chacun exécutant ce programme avec sa sortie redirigée vers \texttt{/dev/null}.
Après avoir lancé toutes les copies, il attend la fin de chaque fils,
ne laissant aucun zombie derrière lui.

\begin{lstlisting}[caption={Constante définissant le nombre de copies}]
#define COPIES 1000000
\end{lstlisting}

\subsection{Comportement face aux limites du système}

En pratique, \texttt{pid\_max} du noyau et la limite utilisateur
\texttt{RLIMIT\_NPROC} sont atteintes bien avant qu'un million de processus
soient créés. Lorsque \texttt{fork} échoue, il retourne \texttt{-1} ;
\texttt{spam} affiche le message d'erreur via \texttt{perror}, passe à
l'itération suivante et continue, de sorte que le compteur \texttt{launched}
reflète fidèlement le nombre de processus réellement créés.

\subsection{Nettoyage}

Toutes les sorties des fils sont redirigées vers \texttt{/dev/null} (ouvert
une seule fois avant les boucles et partagé entre tous les fils). Après tous
les forks, le parent entre dans une boucle bloquante
\texttt{waitpid(-1, NULL, 0)} qui tourne jusqu'à ce qu'il n'y ait plus de
fils, garantissant l'absence de zombies dans la table des processus.

\subsection{Utilisation}

\begin{lstlisting}[language=bash, style=cstyle]
./spam yes          # forker 1 000 000 copies de 'yes'
./spam ls cat       # inonder avec deux programmes en sequence
\end{lstlisting}

\textbf{Avertissement :} ce programme épuise intentionnellement les ressources
système et ne doit être exécuté que dans un environnement contrôlé ou dans un
conteneur avec des paramètres \texttt{ulimit} appropriés.

\subsection{Code source}

\lstinputlisting[caption={\texttt{spam.c} --- utilitaire d'inondation de processus}]{spam.c}

% ── Système de compilation ────────────────────────────────────────────────────
\section{Système de compilation}

Les trois programmes sont compilés par le \texttt{Makefile} suivant. La règle
de motif \texttt{\%: \%.c} compile chaque fichier \texttt{.c} trouvé dans le
répertoire en un binaire portant le même nom de base, en utilisant \texttt{gcc}
avec \texttt{-Wall -Wextra} pour activer des avertissements complets.

\lstinputlisting[language=make, caption={Makefile}]{Makefile}

La commande \texttt{make} compile les trois binaires ; \texttt{make clean}
les supprime.

% ── Bilan ─────────────────────────────────────────────────────────────────────
\section{Bilan comparatif}

\begin{center}
\begin{tabular}{@{}llll@{}}
\toprule
Programme  & API principale                         & Processus créés          & Nettoyage \\
\midrule
myls       & \texttt{opendir}, \texttt{lstat}       & 0 (mono-processus)       & sans objet \\
oroboros   & \texttt{fork}, \texttt{execvp}, \texttt{kill} & un par argument    & SIGTERM + SIGKILL \\
spam       & \texttt{fork}, \texttt{execvp}         & jusqu'à $10^6$ par argument & boucle \texttt{waitpid} \\
\bottomrule
\end{tabular}
\end{center}

Les trois programmes forment une progression naturelle en complexité :
\texttt{myls} opère entièrement dans un seul processus et explore l'API du
système de fichiers ; \texttt{oroboros} introduit la création de processus,
l'exécution et la terminaison contrôlée via des signaux ; \texttt{spam} pousse
les mécanismes de création de processus à leurs limites, servant à la fois de
test de stress et de démonstration des contraintes de la table des processus
du noyau.

\end{document}
